\section{Security Analysis}
\label{sec:security}

We analyze 6FLAS against the threat model of \cref{sec:threat}. Throughout, let
$\skey_\domid$ be the secret domain key, assume the keyed PRF/MAC is a secure
pseudorandom function, and recall that the Flow Label is immutable in transit.

\subsection{Resilience to Known Attacks}
\label{sec:attacks}

We examine 6FLAS against the four attack classes that follow from the
threat model: forgery and tampering, replay, man-in-the-middle
manipulation, and stateless denial of service. 

\paragraph{Forgery and tampering (R1, R2).}
A packet is accepted only if the carried tag equals
$\mtag=\mac_{\skey_\domid}(\fl\,\|\,\textsf{SID}\,\|\,\domid\,\|\,\seq)|_{32}$
(\cref{eq:tag}, \cref{alg:verify}). An adversary without $\skey_\domid$ who
modifies any covered field, the flow identity $\fl$, the path policy
\textsf{SID}, the domain id, or the sequence number, must produce a matching
tag.

\begin{proposition}[Forgery bound]
\label{prop:forgery}
Against an adversary making $q$ verification queries and lacking $\skey_\domid$,
the probability that any single forged or tampered packet is accepted is at most
$2^{-32}+\epsilon_{\prf}$, where $\epsilon_{\prf}$ is the PRF-distinguishing
advantage for the keyed MAC.
\end{proposition}

The $2^{-32}$ term is the chance of guessing the 32-bit truncated tag; a secure
MAC contributes only $\epsilon_{\prf}$. We prove this in \cref{sec:formal}.
Truncating to 32 bits trades tag length for header
compactness; a wider tag can be used when a stronger bound is required. Because the tag covers \textsf{SID}, an adversary also cannot
redirect a packet onto an unauthorized path without invalidating it (R4).

\paragraph{Replay (R3) — probabilistic guarantee.}
Each accepted packet advances the per-direction sliding window $W_v$
(\cref{alg:verify}, lines 4--6, 12). Because the window is shared among
\emph{all flows in a direction} rather than keyed per flow, R3 carries a
probabilistic qualifier: replay of $(\fl_A, \seq_n)$ is detected while
$\seq_n$ remains within the $w$-entry window, but escapes only if other flows
in the same direction collectively advance the window base past $\seq_n$ within
the timing tolerance $\Delta$.

\begin{proposition}[Replay-window eviction bound]
\label{prop:replay}
Let $k = |F_{d,\domid}|$ be the number of flows sharing direction $d$ and
domain $\domid$ in a single replay-window bucket, and let $R_{d,\domid}$ be
their aggregate packet rate. A replayed $\seq_n$ is guaranteed rejected for
at least $\lfloor w/R_{d,\domid} \rfloor$ seconds after first acceptance
($w=256$); at $R_{d,\domid}=10^4$\,pkt/s this covers $25$\,ms. Beyond that
interval, clock-based rejection (tolerance $\Delta$) applies. The probability
of false-positive rejection from two distinct flows coinciding on the same
current sequence number is $O(k^2/2^{24})$, valid when
$k \ll \sqrt{2^{24}} \approx 4{,}096$. With $256$ domains and $6$ forwarding
directions the worst-case bucket at $10^6$ total flows holds
$\lfloor 10^6/(256{\times}6) \rfloor \approx 651$ flows, well within this bound.
\end{proposition}

At the operating rates of \cref{sec:results}, each per-direction window retains
a sequence number for several milliseconds; with $\Delta$ set to cover
propagation delay, the combined guarantee is sufficient for the LEO threat model.
Clock-based rejection is the primary mechanism for long-lived flows whose
sequence numbers have long since left the sliding window.

\paragraph{Man-in-the-middle.}
An on-path adversary can observe and relay packets but cannot alter the covered
fields without breaking \cref{prop:forgery}, and cannot splice a valid tag from
one flow onto another because the tag binds $\fl$ (the per-flow identity) to the
SID and sequence number. Downgrade to plain SRv6 is detected because a 6FLAS
domain expects the \texttt{End.SEC} function and a valid tag; a packet lacking
them is dropped at the first 6FLAS hop.

\paragraph{Denial of service.}
Stateful protocols are vulnerable to state-exhaustion DoS: an attacker opens
many bogus flows to fill the SA table. 6FLAS keeps no per-flow state, so this
vector is closed by construction. The first action in \cref{alg:verify} is a
domain-id check, and verification is a single MAC plus an $O(1)$ window probe,
so the per-packet work an attacker can induce is bounded and small. A Flow-Label
allow-list per domain further lets a node reject traffic outside its assigned
label ranges before computing any MAC.

Taken together, the four analyses confirm that 6FLAS satisfies R1--R4 under the
threat model. Forgery and tampering are bounded by the MAC truncation
(\cref{prop:forgery}); replay is probabilistically rejected by the sliding window (\cref{prop:replay}); man-in-the-middle
manipulation is detected because the tag binds flow identity, path policy, and
sequence number jointly; and denial-of-service via state exhaustion is
structurally impossible because no per-flow state exists to exhaust. The cost of
each check is $O(1)$ in the flow count, consistent with the feasibility
condition of \cref{sec:formulation}.

\subsection{Comparative Security Properties}
\label{sec:sec-compare}

\cref{tab:secprop} compares the security properties 6FLAS provides against the
baselines. 6FLAS matches IPsec/DTLS on data-origin authentication and integrity,
provides a probabilistic anti-replay guarantee (shared per-direction window
rather than per-flow, see \cref{prop:replay}), and uniquely offers stateless
verification and topology-aware access control; it deliberately omits payload
confidentiality, which it leaves to a composable encryption layer.

\begin{table}[t]
\centering
\caption{Security properties provided ($\bullet$ = yes, $\circ$ = partial, --
= no/not applicable).}
\label{tab:secprop}
\renewcommand{\arraystretch}{1.2}
\setlength{\tabcolsep}{4pt}
\begin{tabular}{lcccc}
\toprule
\textbf{Property} & \textbf{IPsec} & \textbf{DTLS} & \textbf{SRv6-Sec (HMAC)} & \textbf{6FLAS} \\
\midrule
Data-origin auth.\ (R1) & $\bullet$ & $\bullet$ & $\bullet$ & $\bullet$ \\
Integrity (R2)          & $\bullet$ & $\bullet$ & $\bullet$ & $\bullet$ \\
Anti-replay (R3)        & $\bullet$ & $\bullet$ & $\circ$   & $\circ$ \\
Path access control (R4)& $\circ$   & --        & $\circ$   & $\bullet$ \\
Stateless verification  & --        & --        & --        & $\bullet$ \\
Handover w/o re-key      & --        & --        & $\circ$   & $\bullet$ \\
Payload confidentiality & $\bullet$ & $\bullet$ & --        & -- \\
\bottomrule
\end{tabular}
\end{table}

\subsection{Formal Security Argument}
\label{sec:formal}

The authentication guarantee is made precise by reducing it to the
unforgeability of the keyed MAC.

\begin{proof}[Proof of \cref{prop:forgery}]
Let $\textsf{MAC}_{\skey}$ be the keyed MAC of \cref{eq:tag} and let
$\textsf{MAC}_{\skey}|_{32}$ be its $32$-bit truncation. Consider an adversary
$\mathcal{A}$ that, after observing honestly generated packets, submits a packet
whose covered message $M=(\fl\,\|\,\textsf{SID}\,\|\,\domid\,\|\,\seq)$ was never
produced by an honest source. (If $M$ equals an honest message, the packet is a
replay and is rejected by the sliding window of \cref{alg:verify}, so we may
assume $M$ is fresh.) Acceptance requires the carried tag to equal
$\textsf{MAC}_{\skey}(M)|_{32}$.

Build a forger $\mathcal{B}$ against $\textsf{MAC}_{\skey}$: $\mathcal{B}$
forwards $\mathcal{A}$'s observation requests to its own tagging oracle to
generate honest packets, and outputs $\mathcal{A}$'s fresh pair
$(M,\mtag)$ as a forgery. Because $M$ was never queried, this is a valid forgery
attempt. If $\textsf{MAC}_{\skey}$ is a secure PRF with distinguishing advantage
$\epsilon_{\prf}$, the probability that $\mtag$ matches the true full tag is at
most $\epsilon_{\prf}$ beyond random guessing. With $32$-bit truncation, an
adversary that simply guesses succeeds with probability $2^{-32}$. Combining,
the per-packet acceptance probability of a fresh forgery is at most
$2^{-32}+\epsilon_{\prf}$. Summing over $q$ verification attempts gives an
overall bound of $q\,(2^{-32}+\epsilon_{\prf})$, the standard multi-query
restatement of the single-packet bound.
\end{proof}

The reduction shows the structure of the guarantee: a successful network
adversary $\mathcal{A}$ that produces an accepted packet it never observed from
an honest source is exactly a MAC forger, because the covered fields it submits
differ from any honestly produced message (otherwise the packet is a replay,
caught by the window). Hence 6FLAS's authentication security reduces to the
unforgeability of its keyed MAC, with the $32$-bit truncation contributing the
$2^{-32}$ term of \cref{prop:forgery}.
